\documentclass[11pt,a4paper]{article}
\usepackage{gtu-el-yazisi}
\begin{document}
\notbilgi{İşaretler ve Sistemler}{Uygulama -- Kararlılık ve konvolüsyon soruları}{salihahocauyg.pdf}
% \notbilgi already prints the bilingual "transcribed with Claude, may contain errors" box.

\sayfa{1}
\begin{ornek}[8)]
Kararlılık: \ $\displaystyle\int_{-\infty}^{\infty} |h(\tau)|\,d\tau < \infty$
\end{ornek}
\[
h(t) = 3^t u(-t+10) \longrightarrow \int_{-\infty}^{10} 3^t\,dt = \frac{3^t}{\log 3}\Big|_{-\infty}^{10}
 = \frac{1}{\log 3}\left(3^{10} - 3^{-\infty}\right) = \underset{\text{Kararlı}}{3^{10}/\log 3}
\]
\vspace{-2mm}\hfill$\longleftarrow\ \log(x) = \ln(x)$\ \ {\small Natural logarithm}

\[
h(t) = e^{-(1-2j)t}u(t) \longrightarrow \int_0^{\infty} e^{-(1-2j)t}\,dt = \frac{e^{-(1-2j)t}}{-(1-2j)}\Big|_0^{\infty}
 = \frac{-1}{1-2j}\left(e^{-\infty} - e^{-0}\right) = \underset{\text{Kararlı}}{\frac{1}{1-2j}}
\]
\[
h(t) = e^{-t}\cos(2t)u(t) \longrightarrow \int_0^{\infty} e^{-t}\cos(2t)\,dt \underset{\text{Kararlı}}{< \infty}
\]
\[
h(t) = t\cos\left(\frac{\pi t}{4}\right)u(t) \longrightarrow \int_0^{\infty} t\cos\left(\frac{\pi t}{4}\right) \longrightarrow \underset{\text{Kararsız}}{\infty}
\]

\begin{ornek}[9)]
$x[n] = u[n-3]$ \ ve \ $h[n] = -\delta[n] + \delta[n+2]$ \ ise \ $y[n] = ?$
\end{ornek}
\[
\begin{aligned}
y[n] = x[n] * h[n] &= u[n-3] * \bigl(-\delta[n] + \delta[n+2]\bigr) = -\delta[n] * u[n-3] + \delta[n+2] * u[n-3]\\
&= -u[n-3] + u[n-1] = \delta[n-1] + \delta[n-2]
\end{aligned}
\]

\sayfa{2}
2) \vurgu[pembe]{Herhangi bir sonlu $t$ aralığında},
\begin{itemize}[label=--]
  \item sonlu sayıda \vurgu[pembe]{maksimum ve minimum} değere sahip olmalıdır.
  \item sonlu sayıda \vurgu[pembe]{süreksizliğe} ve bu noktalarda sonlu değere sahip olmalıdır.
\end{itemize}
$\Rightarrow$ Frekans spektrumları:

\sayfa{3}
{\small\itshape (Claude notu: sayfanın üstü kesik; önceki sorudan yalnızca şu kısım görünüyor:)} \quad
$\underset{t=-3}{\downarrow} \qquad \underset{t=1}{\downarrow}$
% \kok{eksen başı}{eksen sonu}{n/h kökleri}{sıfır noktaları}{eksen etiketleri}{ek çizim}{alt yazı}
\newcommand{\kok}[7]{%
\begin{tikzpicture}[x=0.42cm,y=0.75cm,baseline={(0,0.3)}]
  \draw[-{Stealth},murekkep] (#1,0) -- (#2,0) node[right,font=\small] {$n$};
  \foreach \n/\h in {#3} {\draw[murekkep,thick] (\n,0) -- (\n,\h); \fill[murekkep] (\n,\h) circle (1.5pt);}
  \foreach \n in {#4} {\fill[murekkep] (\n,0) circle (1.5pt);}
  \foreach \n in {#5} {\node[below,font=\scriptsize,inner sep=1.5pt] at (\n,0) {$\n$};}
  #6
  \node[font=\small] at ({(#1+#2)/2},-0.85) {#7};
\end{tikzpicture}}

\begin{ornek}[4)]
$u[n-1] - u[n-2] = ?$

\medskip
\noindent
\kok{-0.6}{5.6}{1/1,2/1,3/1,4/1}{0}{0,1,2,3,4}{\node[font=\scriptsize,left] at (1,1) {$1$}; \node at (4.9,0.45) {$\cdots$};}{$u[n-1]$}
\ $-$\
\kok{-0.6}{5.6}{2/1,3/1,4/1}{0,1}{0,1,2,3,4}{\node[font=\scriptsize,left] at (2,1) {$1$}; \node at (4.9,0.45) {$\cdots$};}{$u[n-2]$}
\ $=$\
\kok{-0.6}{4.6}{1/1}{0,2,3}{0,1,2,3}{\node[font=\scriptsize,left] at (1,1) {$1$}; \node at (4,0.45) {$\cdots$};}{}
\ $= \delta[n-1]$
\end{ornek}

\begin{ornek}[5)]
$u[n-7] - u[n-5] = ?$

\medskip
\noindent
\kok{5.4}{11.6}{7/1,8/1,9/1,10/1}{}{7,8,9,10}{\node[font=\scriptsize,left] at (7,1) {$1$}; \node at (11,0.45) {$\cdots$};}{$u[n-7]$}
\ $-$\
\kok{4.4}{10.6}{5/1,6/1,7/1,8/1,9/1}{}{5,6,7,8,9}{\node[font=\scriptsize,left] at (5,1) {$1$}; \node at (10,0.45) {$\cdots$};}{$u[n-5]$}
\ $=$\
\kok{3.6}{7.6}{5/-1,6/-1}{}{}{\node[font=\scriptsize,above] at (5,0) {$5$}; \node[font=\scriptsize,above] at (6,0) {$6$}; \node[font=\scriptsize,left] at (5,-1) {$1$};}{}
\ $= -\delta[n-5] - \delta[n-6]$
\end{ornek}

\sayfa{4}
\newcommand{\iptal}[1]{\tikz[baseline=(gtuIptal.base)]{\node[inner sep=1pt](gtuIptal){$\displaystyle #1$};
  \draw[turuncukalem,line width=0.7pt] (gtuIptal.south west) -- (gtuIptal.north east);}}
{\small\itshape (Claude notu: bu sayfada turuncu kalemle çok sayıda altı çizme / parantez işareti var; sayfanın üstünde turuncu ``$f_0$'' yazılmış.)}

\begin{ornek}[Ex]
$X(f) = \delta(f-f_0)$ \ ise \ $x(t) = ?$
\end{ornek}
\[
x(t) = \int_{-\infty}^{\infty} X(f)\,e^{j2\pi ft}\,df = \int_{-\infty}^{\infty}\delta(f-f_0)\,e^{j2\pi ft}\,df
 = \int_{-\infty}^{\infty}\delta(f-f_0)\,e^{j2\pi f_0 t}\,df
 = e^{j2\pi f_0 t}\underbrace{\int_{-\infty}^{\infty}\delta(f-f_0)\,df}_{=1} = \underline{\underline{e^{j2\pi f_0 t}}}
\]

\begin{ornek}[Ex]
$X(f) = \delta(f)$ \ ise \ $x(t) = ?$
\end{ornek}
\[
x(t) = \int_{-\infty}^{\infty} X(f)\,e^{j2\pi ft}\,df = \int_{-\infty}^{\infty}\delta(f)\,e^{j2\pi ft}\,df
 = \int_{-\infty}^{\infty}\delta(f)\,\iptal{e^{j2\pi 0 t}}\,df = \int_{-\infty}^{\infty}\delta(f)\,df = 1
\]
\[
\left.\begin{aligned}
\delta(t) &\xleftrightarrow{\ \mathcal{F}\ } 1\\
1 &\xleftrightarrow{\ \mathcal{F}\ } \delta(f)
\end{aligned}\ \right)
\]
\[
\begin{aligned}
e^{j2\pi f_0 t} &\xleftrightarrow{\ \mathcal{F}\ } \underline{\delta(f-f_0)}\\
e^{-j2\pi f_0 t} &\xleftrightarrow{\ \mathcal{F}\ } \delta(f+f_0)
\end{aligned}
\]

\sayfa{5}
{\small\itshape (Claude notu: bu sayfa, el yazısı notun bir PDF görüntüleyicideki ekran görüntüsü; üstteki araç çubuğu aktarılmadı. s.~3, bu sayfanın devamıdır.)}

\begin{ornek}[1)]
\[
\sqrt{2}\,e^{j9\pi/4} = \sqrt{2}\left(\cos 9\pi/4 + j\sin 9\pi/4\right) = \sqrt{2}\left(\cos\pi/4 + j\sin\pi/4\right)
 = \sqrt{2}\left(\frac{\sqrt{2}}{2} + j\frac{\sqrt{2}}{2}\right) = 1+j
\]
\end{ornek}

\begin{ornek}[2)]
$x(t) = u(t) + u(-t)$

\medskip
\noindent
\begin{tikzpicture}[x=0.5cm,y=0.55cm,baseline=(current bounding box.center)]
  \draw[-{Stealth},murekkep] (-2.2,0) -- (2.8,0) node[right,font=\small] {$t$};
  \draw[-{Stealth},murekkep] (0,-0.8) -- (0,1.8) node[right,font=\small] {$u(t)$};
  \draw[murekkep,thick] (0,1) -- (2.2,1);
  \node[left,font=\scriptsize] at (0,1) {$1$};
\end{tikzpicture}\hfill
\begin{tikzpicture}[x=0.5cm,y=0.55cm,baseline=(current bounding box.center)]
  \draw[-{Stealth},murekkep] (-2.2,0) -- (2.8,0) node[right,font=\small] {$t$};
  \draw[-{Stealth},murekkep] (0,-0.8) -- (0,1.8) node[right,font=\small] {$u(-t)$};
  \draw[murekkep,thick] (-2.2,1) -- (0,1);
  \node[right,font=\scriptsize] at (0,1) {$1$};
\end{tikzpicture}\hfill
\begin{tikzpicture}[x=0.5cm,y=0.55cm,baseline=(current bounding box.center)]
  \draw[murekkep] (-2.2,0) -- (2.8,0);
  \draw[-{Stealth},murekkep] (0,-0.8) -- (0,1.8) node[right,font=\small] {$x(t)$};
  \draw[murekkep,thick] (-2.2,1) -- (2.8,1);
\end{tikzpicture}\hfill
\begin{minipage}[c]{0.27\linewidth}
$T \to \infty$\\
Periyodik değildir
\end{minipage}
\end{ornek}

\begin{ornek}[3)]
$x(t) = 2, \ -1 < t < 3$ \ ise \ $x(1-2t) + x(2+t) = ?$
\end{ornek}

\noindent
\begin{minipage}[t]{0.5\linewidth}
\vspace{0pt}
\begin{tikzpicture}[x=0.42cm,y=0.42cm]
  \draw[-{Stealth},murekkep] (-2.5,0) -- (4.5,0) node[right,font=\small] {$t$};
  \draw[-{Stealth},murekkep] (0,-0.8) -- (0,3) node[right,font=\small] {$x(t)$};
  \draw[murekkep,thick] (-1,0) -- (-1,2) -- (3,2) -- (3,0);
  \node[above right,font=\scriptsize,inner sep=1pt] at (0,2) {$2$};
  \node[below,font=\scriptsize] at (-1,0) {$-1$};
  \node[below,font=\scriptsize] at (3,0) {$3$};
\end{tikzpicture}
\[
x(1-2t) \longrightarrow x(t) \quad t:[-1,3]
\]
\[
\begin{array}{c}
1-2t : [-1,3]\\
\swarrow \qquad\qquad \searrow\\
\begin{array}{cc} 1-2t = -1 \quad & \quad 1-2t = 3\\ t = 1 \quad & \quad t = -1 \end{array}
\end{array}
\]
\[
x(2+t) \longrightarrow x(t) \quad t:[-1,3]
\]
\[
\begin{array}{c}
2+t : [-1,3]\\
\swarrow \qquad \searrow\\
\begin{array}{cc} t = -3 \quad & \quad t = 1 \end{array}
\end{array}
\]
\end{minipage}\hfill
\begin{minipage}[t]{0.47\linewidth}
\vspace{0pt}
\centering
\begin{tikzpicture}[x=0.42cm,y=0.42cm]
  \draw[-{Stealth},murekkep] (-2.5,0) -- (3,0) node[right,font=\small] {$t$};
  \draw[-{Stealth},murekkep] (0,-0.8) -- (0,3) node[right,font=\small] {$x(1-2t)$};
  \draw[murekkep,thick] (-1,0) -- (-1,2) -- (1,2) -- (1,0);
  \node[above right,font=\scriptsize,inner sep=1pt] at (0,2) {$2$};
  \node[below,font=\scriptsize] at (-1,0) {$-1$};
  \node[below,font=\scriptsize] at (1,0) {$1$};
\end{tikzpicture}

\medskip
\begin{tikzpicture}[x=0.42cm,y=0.42cm]
  \draw[-{Stealth},murekkep] (-4,0) -- (3,0) node[right,font=\small] {$t$};
  \draw[-{Stealth},murekkep] (0,-0.8) -- (0,3) node[right,font=\small] {$x(t+2)$};
  \draw[murekkep,thick] (-3,0) -- (-3,2) -- (1,2) -- (1,0);
  \node[above right,font=\scriptsize,inner sep=1pt] at (0,2) {$2$};
  \node[below,font=\scriptsize] at (-3,0) {$-3$};
  \node[below,font=\scriptsize] at (1,0) {$1$};
\end{tikzpicture}

\medskip
\begin{tikzpicture}[x=0.42cm,y=0.42cm]
  \draw[-{Stealth},murekkep] (-4,0) -- (3,0) node[right,font=\small] {$t$};
  \draw[-{Stealth},murekkep] (0,-0.8) -- (0,5) node[right,font=\small] {$x(1-2t)+x(t+2)$};
  \draw[murekkep,thick] (-3,0) -- (-3,2) -- (-1,2) -- (-1,4) -- (1,4) -- (1,0);
  \draw[murekkep,dotted] (-1,2) -- (0,2);
  \node[right,font=\scriptsize,inner sep=1pt] at (0,2) {$2$};
  \node[above right,font=\scriptsize,inner sep=1pt] at (0,4) {$4$};
  \node[below,font=\scriptsize] at (-3,0) {$-3$};
  \node[below,font=\scriptsize] at (-1,0) {$-1$};
  \node[below,font=\scriptsize] at (1,0) {$1$};
\end{tikzpicture}
\end{minipage}

\medskip
{\small\itshape (Claude notu: sayfanın altındaki 4.\ soru, $u[n-1]-u[n-2] = \delta[n-1]$, s.~3'teki ile aynıdır; tekrar aktarılmadı.)}

\noindent\sayfa{6}
{\small\itshape (Claude notu: bu sayfa s.~5'teki 3.\ sorunun aynısıdır (aynı çizimler ve çözüm); tek fark, soru ifadesindeki $x(t) = 2$, \ $-1 < t < 3$, \ $x(1-2t)$ ve $x(2+t)$ kısımlarının turuncu kalemle altının çizilmiş olmasıdır. Tekrar aktarılmadı.)}

\sayfa{7}
{\small\itshape (Claude notu: sayfanın üstündeki 5.\ soru, $u[n-7]-u[n-5] = -\delta[n-5]-\delta[n-6]$, s.~3'teki ile aynıdır; tekrar aktarılmadı. 6.\ soru turuncu kalemle çerçeve içine alınmış. Sağ alttaki CamScanner logosu aktarılmadı.)}

\begin{ornek}[6)]
$\cos\left(\frac{\pi}{3}n\right)\delta[n+4] = ?$
\end{ornek}
\[
\text{Reminder:}\quad
\begin{aligned}[t]
x[n]\,\delta[n] &= x[0]\,\delta[n]\\
x[n]\,\delta[n-n_0] &= x[n_0]\,\delta[n-n_0]
\end{aligned}
\]
\[
\cos\left(\frac{\pi}{3}n\right)\delta[n+4] = \cos\left(-\frac{4\pi}{3}\right)\delta[n+4] = \cos\left(\frac{4\pi}{3}\right)\delta[n+4]
\]

\end{document}
